%% 47장 비선형 최소제곱
\chapter{Nonlinear Least Squares}
\label{ch:nonlinear-least-squares}

\abstract{모형을 데이터에 맞추는 일의 대부분은 잔차의 제곱합을 최소화하는 문제이다. 이 목적 함수에는 Hessian의 한 부분을 야코비안만으로 계산할 수 있다는 특별한 구조가 있다. 이 장에서는 그 구조를 쓰는 세 방법을 다룬다. Gauss--Newton 방법은 Hessian의 나머지 부분을 버리고 선형 최소제곱 문제를 되풀이해 푼다. 그 수렴 속도가 버린 부분의 크기, 곧 최소점의 잔차에 달려 있음을 오차의 점화식으로 보이고 잔차가 크면 수렴하지 않는 예를 수치로 확인한다. Levenberg--Marquardt 방법은 45장의 신뢰 영역을 최소제곱에 맞춘 것이며 여기서는 MINPACK의 \code{lmder}가 열 피벗 QR과 Givens 회전으로 부분문제를 푸는 과정을 따라간다. 끝으로 변수에 상하한이 있을 때 쓰는 반사 신뢰 영역 방법(TRF)이 Coleman--Li 척도로 경계를 다루는 방식을 설명한다.}

\section{The Structure of Least-Squares Problems}
\label{sec:nls-structure}

잔차 벡터 $r:\mathbb R^n\to\mathbb R^m$($m\ge n$)이 주어졌을 때
\[
F(x)=\frac12\|r(x)\|^2=\frac12\sum_{i=1}^mr_i(x)^2
\]
을 최소화하는 문제를 \term{비선형 최소제곱}{nonlinear least squares}이라 한다. 측정값 $y_i$에 모형 $\phi(t_i;x)$를 맞추면 $r_i(x)=\phi(t_i;x)-y_i$이다. 잡음이 평균 $0$, 같은 분산의 독립 정규 분포를 따르면 이 최소화는 최대우도 추정이 되고 그 통계적 성질은 67장에서 다룬다. $r$의 야코비안을 $J(x)\in\mathbb R^{m\times n}$이라 하면
\[
\nabla F(x)=J(x)^{\mathsf T}r(x),\qquad\nabla^2F(x)=J(x)^{\mathsf T}J(x)+S(x),\qquad S(x)=\sum_{i=1}^mr_i(x)\nabla^2r_i(x)
\]
이다. 기울기와 Hessian의 첫째 항 $J^{\mathsf T}J$는 야코비안만으로 계산된다. 둘째 항 $S$에는 잔차마다 Hessian이 필요하지만 각 항에 잔차 $r_i$가 곱해져 있으므로 모형이 데이터에 잘 맞아 잔차가 작으면 $S$도 작다. 모형이 선형이면 $\nabla^2r_i=0$이어서 $S=0$이고, 문제는 22장의 선형 최소제곱이 된다.

\section{The Gauss--Newton Method}
\label{sec:gauss-newton}

\subsection{Linearization}

Newton 방법의 Hessian에서 $S$를 버리면
\[
J_k^{\mathsf T}J_kp_k=-J_k^{\mathsf T}r_k
\]
가 되는데, 이것은 선형 최소제곱 문제
\[
\min_p\|J_kp+r_k\|
\]
의 정규방정식이다. 잔차를 $r(x_k+p)\approx r_k+J_kp$로 선형화하고 그 선형 문제를 푸는 이 방법이 \term{Gauss--Newton 방법}{Gauss--Newton method}이다. 정규방정식을 직접 풀면 조건수가 제곱되므로 보통은 21장의 QR 분해 $J_k=QR$로 $Rp_k=-Q^{\mathsf T}r_k$를 푼다. $J_k$의 열이 일차 독립이면 $J_k^{\mathsf T}J_k$가 양정치이므로 $p_k$는 하강 방향이다. 준위 집합 근방에서 $J$가 유계이고 그 특이값이 $0$에서 고르게 떨어져 있으면 44장의 직선 탐색과 함께 쓸 때 Zoutendijk의 정리로 $J_k^{\mathsf T}r_k\to0$이다. Hartley는 1961년에 이 방법에 직선 탐색을 붙인 판을 통계학자들에게 알렸다\cite{hartley1961}.

\subsection{Rate of Convergence}

Gauss--Newton 걸음은 Hessian의 일부를 버렸으므로 Newton 방법처럼 이차 수렴하지는 못한다. 얼마나 느려지는지는 다음 명제로 알 수 있다.

\begin{proposition}
\label{prop:gn-rate}
$r$이 $x_\ast$ 근방에서 두 번 연속 미분 가능하고 그 이계 도함수가 Lipschitz 연속이며, $\nabla F(x_\ast)=0$이고 $J_\ast=J(x_\ast)$의 열이 일차 독립이라 하자. 단위 걸음의 Gauss--Newton 반복은 $e_k=x_k-x_\ast$에 대해
\[
e_{k+1}=-(J_\ast^{\mathsf T}J_\ast)^{-1}S(x_\ast)\,e_k+O(\|e_k\|^2)
\]
를 만족한다. 따라서 행렬 $(J_\ast^{\mathsf T}J_\ast)^{-1}S(x_\ast)$의 스펙트럼 반지름이 $1$보다 작으면 $x_\ast$ 근방에서 선형 수렴하고, $S(x_\ast)=0$이면 이차 수렴한다.
\end{proposition}

\begin{proof}
$A_k=J_k^{\mathsf T}J_k$로 쓰자. $\nabla F(x_\ast)=0$이므로 Taylor 정리로 $\nabla F(x_k)=\nabla^2F(x_\ast)e_k+O(\|e_k\|^2)=(A_\ast+S_\ast)e_k+O(\|e_k\|^2)$이다. 걸음 $p_k=-A_k^{-1}\nabla F(x_k)$에서 $A_k^{-1}=A_\ast^{-1}+O(\|e_k\|)$이므로
\begin{align*}
e_{k+1}&=e_k-A_k^{-1}\nabla F(x_k)=e_k-A_\ast^{-1}(A_\ast+S_\ast)e_k+O(\|e_k\|^2)\\
&=-A_\ast^{-1}S_\ast e_k+O(\|e_k\|^2)
\end{align*}
이다.
\end{proof}

$S(x_\ast)$의 크기는 최소점의 잔차와 모형 곡률의 곱이다. 그래서 잔차가 $0$인 문제, 곧 모형이 데이터를 정확히 지나는 문제에서는 Gauss--Newton이 Newton 방법만큼 빠르고 잔차가 작으면 빠르게 선형 수렴하며 잔차가 크면 느리게 수렴하거나 발산한다. Dennis와 Schnabel의 예\cite{dennis1996}로 확인해 보자. 한 변수 모형 $\phi(t;x)=e^{xt}$를 세 점 $(1,2)$, $(2,4)$, $(3,y_3)$에 맞춘다. $y_3=8$이면 $x_\ast=\ln2$에서 잔차가 정확히 $0$이고, $y_3$을 $8$에서 멀리 둘수록 최소점의 잔차가 커진다. \cref{prop:gn-rate}의 비 $-S_\ast/(J_\ast^{\mathsf T}J_\ast)$를 계산하고 $x_\ast+0.1$에서 Gauss--Newton을 돌린 결과가 \cref{tab:gn-rate}이다. 연속한 오차의 비가 예측한 비와 맞고 $y_3=-4$에서는 비가 $-2.2$여서 최소점 바로 옆($x_\ast+10^{-3}$)에서 시작해도 오차가 $-2.2\times10^{-3}$, $4.9\times10^{-3}$, $-1.1\times10^{-2}$, $2.4\times10^{-2}$으로 부호를 바꾸며 커진다.

\begin{table}[t]
\begin{sdtable}{colspec={ccccc}, row{1}={font=\small\bfseries}}
\toprule
$y_3$ & $x_\ast$ & $F(x_\ast)$ & $-S_\ast/(J_\ast^{\mathsf T}J_\ast)$ & Gauss--Newton 오차의 처음 넷 \\
\midrule
$8$ & $0.693147$ & $0$ & $0$ & $1.3\times10^{-2}$, $2.5\times10^{-4}$, $9.0\times10^{-8}$, $1.2\times10^{-14}$ \\
$6$ & $0.610176$ & $0.228$ & $-0.012$ & $1.2\times10^{-2}$, $7.5\times10^{-5}$, $-8.8\times10^{-7}$, $1.0\times10^{-8}$ \\
$3$ & $0.440050$ & $1.64$ & $-0.060$ & $8.6\times10^{-3}$, $-4.0\times10^{-4}$, $2.4\times10^{-5}$, $-1.4\times10^{-6}$ \\
$-1$ & $0.044744$ & $6.98$ & $-0.473$ & $-2.3\times10^{-2}$, $1.2\times10^{-2}$, $-5.5\times10^{-3}$, $2.7\times10^{-3}$ \\
$-4$ & $-0.371929$ & $16.4$ & $-2.20$ & 수렴하지 않음 \\
\bottomrule
\end{sdtable}
\caption{$e^{xt}$를 $(1,2)$, $(2,4)$, $(3,y_3)$에 맞출 때 Gauss--Newton의 수렴. 오차는 $x_\ast+0.1$에서 시작한 $x_k-x_\ast$이다. 비의 부호가 음수이면 오차의 부호가 결국 번갈아 바뀐다.}
\label{tab:gn-rate}
\end{table}

\begin{aside}[Gauss와 소행성 Ceres]
1801년 1월 Piazzi가 발견한 소행성 Ceres는 40일쯤 관측된 뒤 태양 뒤로 숨어 버렸다. 스물네 살의 Gauss는 그 짧은 관측으로 궤도를 계산해 그해 말 Ceres가 다시 나타날 자리를 맞혔고, 그 방법을 1809년의 「천체 운동론」에 최소제곱법으로 정리했다\cite{gauss1809}. 비선형인 궤도 요소를 선형화해 되풀이 고치는 그의 계산이 Gauss--Newton 방법의 원형이다. 같은 무렵 Legendre도 최소제곱법을 발표해(1805)\cite{legendre1805} 둘 사이에 우선권 다툼이 있었다.
\end{aside}

\section{The Levenberg--Marquardt Method}
\label{sec:levenberg-marquardt}

\subsection{Damping and Trust Regions}

Gauss--Newton의 약점은 둘이다. 잔차가 크면 수렴하지 않을 수 있고, $J$의 열이 거의 일차 종속이면 걸음이 터무니없이 커진다. Levenberg(1944)와 Marquardt(1963)는 정규방정식에 대각 항을 더해
\[
(J_k^{\mathsf T}J_k+\lambda D_k^{\mathsf T}D_k)\,p=-J_k^{\mathsf T}r_k
\]
를 풀었다\cite{levenberg1944,marquardt1963}. $D_k$는 대각인 척도 행렬이다. $\lambda$가 $0$이면 Gauss--Newton 걸음이고, $\lambda$가 크면 걸음이 짧아지며 척도를 준 가장 가파른 내리막 방향으로 돈다. 45장의 Moré--Sorensen 정리에 따르면 이 걸음은 이차 모형 $\frac12\|J_kp+r_k\|^2$를 영역 $\|D_kp\|\le\Delta_k$에서 최소화하는 신뢰 영역 부분문제의 해이다. 모형의 Hessian $J_k^{\mathsf T}J_k$가 언제나 양반정치이고 기울기 $J_k^{\mathsf T}r_k$가 $J_k$의 영공간과 직교하므로, $J_k$의 열이 종속이어도 최소 노름 Gauss--Newton 걸음이 해가 되어 어려운 경우를 따로 다룰 일이 없고 \cref{thm:tr-convergence}의 전역 수렴이 그대로 적용된다. 그래서 \term{Levenberg--Marquardt 방법}{Levenberg--Marquardt method}은 $\lambda$를 직접 조절하는 대신 반지름 $\Delta_k$를 조절하는 신뢰 영역 방법으로 구현하는 것이 정석이다. Moré는 1978년에 이 관점으로 MINPACK의 구현을 설계했다\cite{more1978lm,more1980minpack}.

\subsection{MINPACK's lmder}

\code{lmder}의 한 반복은 다음과 같다\cite{more1978lm}.

첫째, 야코비안 $J$를 열 피벗 Householder QR $J\Pi=QR$로 분해하고 $Q^{\mathsf T}r$를 만든다. 척도 $D=\operatorname{diag}(d_j)$는 $J$의 열 노름이고 반복 사이에 줄지 않도록 지금까지의 최댓값을 쓴다. 척도를 이렇게 정하면 변수의 단위를 바꾸어도 같은 걸음이 나온다.

둘째, 부분문제를 \code{lmpar}가 푼다. 먼저 $\lambda=0$의 Gauss--Newton 걸음이 $\|Dp\|\le1.1\Delta$를 만족하면 그것을 쓴다. 아니면 45장처럼 $1/\|Dp(\lambda)\|$의 Newton 반복으로 $\bigl|\|Dp(\lambda)\|-\Delta\bigr|\le0.1\Delta$인 $\lambda$를 찾는다. 반복마다 $\lambda$가 바뀌어도 QR을 처음부터 다시 하지 않는다. 최소제곱 문제
\[
\min_p\Bigl\|\begin{pmatrix}J\\\sqrt\lambda D\end{pmatrix}p+\begin{pmatrix}r\\0\end{pmatrix}\Bigr\|
\]
의 계수 행렬은 이미 구한 $R$ 아래에 대각 행렬 $\sqrt\lambda\,\Pi^{\mathsf T}D\Pi$를 붙인 것이므로, 대각 원소를 하나씩 Givens 회전으로 $R$에 녹여 넣으면 $m$과 무관한 $O(n^3)$의 일로 새 삼각 행렬을 얻는다(\code{qrsolv}). 처음부터 QR을 다시 하면 $O(mn^2)$이 든다. 정규방정식을 만들지 않으므로 조건수가 제곱되지 않는다.

셋째, 비 $\rho=(\|r(x)\|^2-\|r(x+p)\|^2)/(\|r(x)\|^2-\|r(x)+Jp\|^2)$로 반지름을 고친다. $\rho\le1/4$이면 반지름을 $0.1$에서 $0.5$ 사이의 비율로 줄이고(비율은 실제 감소와 방향 도함수로 만든 이차 보간에서 정한다), $\rho>1/4$이면서 $\rho\ge3/4$이거나 $\lambda=0$이면 $\Delta=2\|Dp\|$로 둔다. $\rho\ge10^{-4}$이면 걸음을 받아들인다. 처음 반지름은 $\Delta_0=100\|Dx_0\|$이다($\|Dx_0\|=0$이면 $100$).

멈춤 판정은 셋이다. 실제와 예측의 상대 감소가 모두 $\mathrm{ftol}$ 이하이면 함수값이 수렴한 것이고, $\Delta\le\mathrm{xtol}\cdot\|Dx\|$이면 위치가 수렴한 것이며, 잔차 벡터와 $J$의 각 열 사이 각의 코사인이 모두 $\mathrm{gtol}$ 이하이면 기울기가 $0$인 것이다. 마지막 판정은 $\|J^{\mathsf T}r\|$을 직접 재지 않고 각으로 재므로 척도에 무관하다. scipy와 stardust의 \code{least\_squares}는 세 허용오차의 기본값이 모두 $10^{-8}$이다.

\cref{tab:gn-rate}의 마지막 줄, Gauss--Newton이 수렴하지 않은 $y_3=-4$의 문제를 Levenberg--Marquardt로 풀면 $x_\ast+10^{-3}$에서 기본 허용오차로 함수값 $4$번 만에 수렴한다. $y_3=-8$에서는 비가 $-6.5$여서 Gauss--Newton의 오차가 세 걸음 만에 $10^{-1}$을 넘지만 Levenberg--Marquardt는 역시 $4$번으로 수렴한다. 잔차가 커서 $S$가 중요한 문제에서도 Levenberg--Marquardt는 수렴은 하되 선형으로만 수렴한다. 그런 문제에는 $S$를 준Newton 방법으로 근사해 더하는 NL2SOL 같은 방법이 있다\cite{dennis1981nl2sol}.

\section{Bounds and Reflective Trust Regions}
\label{sec:trf}

\subsection{Coleman--Li Scaling}

변수에 상하한 $l\le x\le u$가 있으면 반복점이 상자 안에 머물러야 한다. 44장의 L-BFGS-B는 활성 집합을 정하는 방식이었다. Coleman과 Li는 다른 길을 택했다\cite{coleman1996}. $g=\nabla F$라 하면 경계 제약 최소화의 일계 조건은 성분마다 $l_i<x_i<u_i$이면 $g_i=0$, $x_i=l_i$이면 $g_i\ge0$, $x_i=u_i$이면 $g_i\le0$이다. 이 조건은
\[
v_i=\begin{cases}u_i-x_i,&g_i<0,\ u_i<\infty,\\x_i-l_i,&g_i>0,\ l_i>-\infty,\\1,&\text{그 밖}\end{cases}
\]
로 정한 벡터 $v$에 대해 방정식 $\operatorname{diag}(v)\,g=0$과 같다. 내리막 방향 $-g$가 향하는 경계까지의 거리를 곱하므로 그 경계에 닿은 성분은 기울기가 $0$이 아니어도 된다. Coleman--Li 방법은 이 방정식에 Newton 방법을 쓰는 것으로 볼 수 있으며 $x=D\hat x$, $D=\operatorname{diag}(v^{1/2})$로 척도를 바꾼 공간에서 신뢰 영역 걸음을 찾는다. 경계 가까이에서 내리막이 그 경계를 향하는 성분은 $v_i$가 작아 원래 공간의 걸음 성분 $p_i=v_i^{1/2}\hat p_i$가 짧아지므로, 걸음이 경계에 부딪히기 전에 저절로 줄어든다.

\subsection{Reflection and Strictly Interior Iterates}

척도를 바꾼 부분문제는 Hessian 근사 $D J^{\mathsf T}JD$에 경계에서 오는 대각 항을 더한 것이고 scipy의 \texttt{trf}는 이 문제를 SVD로 분해한 뒤 Moré의 방법으로 $\bigl|\|\hat p\|-\Delta\bigr|\le0.01\Delta$까지 푼다. 걸음이 상자를 벗어나면 세 후보 가운데 모형을 가장 많이 줄이는 것을 고른다. 하나는 걸음을 경계까지 자르되 $\theta=\max(0.995,\,1-\|\operatorname{diag}(v)\,g\|_\infty)$배로 조금 덜 가서 엄격한 내부에 머무는 것이다. 다른 하나는 처음 닿은 경계에서 걸음을 거울처럼 \term{반사}{reflection}시키고, 반사된 방향을 따라 상자와 신뢰 영역 안에서 모형을 가장 많이 줄이는 길이만큼 가는 것이다. 마지막은 척도를 바꾼 가장 가파른 내리막 방향으로 신뢰 영역과 상자 안에서 최소화하는 것이다. 반사는 경계에 부딪힌 걸음을 버리지 않고 살려 쓰는 방법이어서 활성 집합을 바꾸는 결정 없이 경계를 따라 미끄러지게 해 준다. 이 방법이 \term{반사 신뢰 영역}{trust region reflective, TRF} 방법이다\cite{branch1999}. 반복점은 언제나 상자의 엄격한 내부에 있고, 최소점이 경계에 있으면 경계에 한없이 다가간다.

Rosenbrock 함수를 잔차 $r=(10(x_2-x_1^2),\,1-x_1)$으로 쓰면 잔차가 $0$인 최소제곱 문제가 된다. $(-1.2,1)$에서 scipy의 \code{'lm'}은 함수값 $21$번, \code{'trf'}는 $25$번으로 최소점 $(1,1)$을 정확히 찾는다. $J^{\mathsf T}J$가 Hessian을 잘 근사하므로 함수값 $39$번이 든 44장의 BFGS보다 적다. 경계 $x_1\le0.8$을 주면 \code{'trf'}는 $30$번으로 $(0.8,0.64)$에 이르고 $x_1$이 위 경계에서 활성이라고 알린다. 그곳의 값은 $F=(1-0.8)^2/2=0.02$이다.

\begin{instardust}[\code{docs/design/optimize.md} §3.4]
\code{sd.optimize.least\_squares}는 \code{'lm'}(MINPACK \code{lmder})과 \code{'trf'}(scipy 1.18.1의 \texttt{trf}, \code{'exact'} 부분문제)를 둔다. 커널(\code{csrc/optimize/lsq.cpp})은 회귀 영역의 \code{sd.regression.nls}와 같으며 \code{lmder}는 한 소스를 두 번 빌드한다. scipy는 MINPACK의 C 판을 clang의 기본 설정으로 빌드해 \code{a*b+c}를 FMA 하나로 계산하므로, \code{least\_squares('lm')}과 \code{root('lm')}은 같은 축약을 켠 빌드로 scipy와 비트까지 같은 걸음을 가고, \code{nls}는 축약 없는 빌드로 자기 골든 출력을 지킨다. $m\ge10n$이고 $mn\ge2000$인 긴 문제에서는 먼저 LAPACK의 Householder QR $J=Q_0R_0$을 한 번 하고, 그 뒤의 열 피벗 QR(\code{lmder})이나 SVD(\code{'trf'})는 $n\times n$ 행렬 $R_0$에 한다. 직교 변환은 열의 노름과 내적을 지키므로 결과는 같고, 야코비안마다 $m$에 비례하는 일은 QR 하나뿐이다. 시험점의 잔차가 유한하지 않으면 그 걸음을 거절하고 반지름을 줄인다. \code{nfev}는 차분 야코비안을 위한 호출까지 센다(scipy는 세지 않는다).
\end{instardust}

\section{Concluding Remarks}

비선형 최소제곱은 목적 함수의 구조가 알고리즘을 정하는 좋은 예이다. 기울기와 Hessian의 한 부분이 야코비안만으로 계산되므로 이차 정보를 거의 공짜로 얻는다. Gauss--Newton은 그 부분만 쓰며 잔차가 작을수록 Newton 방법에 가까워진다. 잔차가 크면 버린 부분이 커져 느려지거나 발산하는데 Levenberg--Marquardt는 45장의 신뢰 영역으로 그 위험을 막는다. MINPACK의 구현은 열 피벗 QR과 Givens 회전으로 $\lambda$가 바뀔 때마다 드는 일을 $m$과 무관한 $O(n^3)$으로 줄이고, 척도와 각으로 잰 판정으로 변수의 단위에 무관하게 동작한다. 경계가 있으면 Coleman--Li 척도와 반사가 활성 집합 없이 경계를 다룬다. 다음 장에서는 경계보다 일반적인 제약과 쌍대성을 다룬다.

비선형 최소제곱의 표준 참고서는 Dennis와 Schnabel의 책\cite{dennis1996} 10장과 Nocedal과 Wright의 책\cite{nocedal2006} 10장이다.

\chapterreferences
